\heading{实验物理中的统计方法-第3章 蒙特卡罗方法}


\begin{question}
写一段小程序，利用 ROOT 中的随机数产生子产生 \(10^4\) 个 \((0,1]\)
之间均匀分布的随机数，将结果画到直方图中（分 100 个区间）。
\begin{solution}
可按如下步骤实现：产生 10000 个均匀随机数，每个随机数填入区间为
\([0,1]\)、分箱数为 100 的直方图。ROOT 中可用
\(\texttt{gRandom->Rndm()}\) 或 \(\texttt{TRandom3::Uniform(0,1)}\) 产生
\([0,1)\) 上的均匀随机数；端点取 \((0,1]\) 或 \([0,1)\)
对连续分布没有实际影响。
\pyfig[0.76\linewidth]{figures/fig_03_01_uniform.png}{Python 复现实验生成的均匀分布直方图；虚线为每箱期望计数 100。}
理论上
\[
E[X]=\int_0^1 x\,dx=\frac12,
\qquad
V[X]=\int_0^1 x^2\,dx-\left(\frac12\right)^2=\frac1{12}.
\]
一次典型模拟可能得到
\[
\bar x\simeq0.497,
\qquad
s^2\simeq0.0830,
\]
与理论值 \(0.5\) 和 \(1/12=0.08333\) 相符；由于样本量有限，具体数值会随随机数种子波动。
\end{solution}
\end{question}


\begin{question}
修改本章第 1 题中的直方图，使其只有 5 个区间，事例数
\(N=100\)。产生的直方图可以看做对矢量 \((n_1,\ldots,n_5)\)
的观测，该矢量满足参数 \(N=100\)，\(p_i=0.2\ (i=1,\ldots,5)\)
的多项分布。

\begin{enumerate}
\item
  将上面的程序代码放到循环中产生直方图，重复该 MC 实验 100
  次，每次用不同的随机数种子。只要保证每次实验程序不中断，程序在下一次实验（即下一个循环）时自动更新种子。定义一个直方图，将每次实验第
  \(i\) 个区间（比如 \(i=3\)）的值填充其中。这应该服从均值为
  \(Np_i=20\)，标准偏差为 \(\sqrt{Np_i(1-p_i)}=4\) 的二项分布。
\item
  为任意两个区间的值 \(n_i\) 和 \(n_j\)
  作出散点图（二维直方图）。其协方差的理论值为
  \(\operatorname{cov}[n_i,n_j]=-Np_ip_j=-4\)，或者说关联系数
  \(\rho=-4/4^2=-0.25\)。
\end{enumerate}
\begin{solution}
每次实验中把 \(N=100\) 个均匀随机数填入 5 个等宽箱。由于每个箱的概率为
\(p_i=0.2\)，单个箱计数满足边缘二项分布
\[
n_i\sim\operatorname{Binom}(100,0.2),
\qquad
E[n_i]=20,
\qquad
\sigma_{n_i}=\sqrt{100\cdot0.2\cdot0.8}=4.
\]
任意两个不同箱的计数满足多项分布协方差
\[
\operatorname{cov}(n_i,n_j)=-Np_ip_j=-100\cdot0.2\cdot0.2=-4,
\qquad
\rho_{ij}=\frac{-4}{4\cdot4}=-0.25.
\]

因此应把 100 次重复实验得到的某个箱计数（如 \(n_3\)）再填入一个新直方图，
其中心应在 20 附近，宽度约为 4。对 \((n_i,n_j)\) 作二维散点图时，点云应表现出负相关：一个箱中的事件数偏多时，其它箱的总数被固定条件压低。

若只重复 100 次，样本协方差和相关系数的统计波动仍会较大；例如得到
\(\operatorname{corr}(n_1,n_2)=-0.14\) 并不奇怪。若把重复次数提高到
\(10^4\) 或更多，样本均值、标准差和相关系数会更稳定地趋近于上面的理论值。
\pyfig[0.68\linewidth]{figures/fig_03_02_bin_count.png}{\(10^4\) 次重复实验中单个箱计数的经验分布及二项分布理论曲线。}
\pyfig[0.54\linewidth]{figures/fig_03_02_scatter.png}{两个不同箱计数的散点图；总数固定导致负相关。}
\end{solution}
\end{question}


\begin{question}
考虑锯齿分布

\[
f(x)=\begin{cases}
\dfrac{2x}{x_{\max}^2}, & 0<x<x_{\max},\\
0, & \text{其它}.
\end{cases}
\]

\begin{enumerate}
\item
  参考课本 3.2 节，利用变换方法寻找函数 \(x(r)\)，以产生服从 \(f(x)\)
  的随机数。用程序实现该方法，并生成一个直方图。（可以取
  \(x_{\max}=1\)。）
\item
  参考课本 3.3
  节，编写一段程序，利用舍选法产生服从锯齿分布的随机数。画出相应的直方图。
\end{enumerate}
\begin{solution}
题设密度为
\[
f(x)=\frac{2x}{x_{\max}^2},\qquad 0<x<x_{\max}.
\]

\begin{enumerate}
\item 反变换法

分布函数为
\[
F(x)=\int_0^x\frac{2t}{x_{\max}^2}\,dt=\frac{x^2}{x_{\max}^2}.
\]
令 \(r\sim U(0,1)\)，取 \(r=F(x)\)，得到
\[
x(r)=x_{\max}\sqrt r.
\]
因此程序中每产生一个均匀随机数 \(r\)，填入
\(x=x_{\max}\sqrt r\) 即可。直方图应随 \(x\) 线性上升。

\item 舍选法

取建议分布
\[
g(x)=\frac1{x_{\max}},\qquad 0<x<x_{\max}.
\]
则
\[
\frac{f(x)}{g(x)}=\frac{2x}{x_{\max}}\le 2,
\]
可取 \(M=2\)。算法为：先在 \((0,x_{\max})\) 上均匀产生候选值
\(x\)，再产生 \(u\sim U(0,1)\)，若
\[
u<\frac{f(x)}{Mg(x)}=\frac{x}{x_{\max}}
\]
则接受该 \(x\)，否则重抽。接受率为 \(1/M=1/2\)。
\end{enumerate}
\pyfig[0.72\linewidth]{figures/fig_03_03_sawtooth.png}{反变换法和舍选法产生的样本直方图；实线为理论密度 \(f(x)=2x\)。}
\end{solution}
\end{question}


\begin{question}
该练习的目的是产生服从高斯分布的随机数。有很多算法可以产生高斯分布，一个最简单的适合教学目的的算法基于中心极限定理：当
\(n\) 很大时，\(n\)
个随机数之和趋向于高斯分布，只要其中任何一项不占绝对份额（参考课本第 10
章）。

\begin{enumerate}
\item
  假设 \(x\) 均匀分布于 \([0,1]\)，考虑 \(n\) 个独立随机变量 \(x\)
  之和，
\end{enumerate}

\[
y=\sum_{i=1}^{n}x_i.
\]

证明 \(y\) 的期待值为 \(n/2\)，方差为 \(n/12\)。并证明变量 \(z\)

\[
z=\frac{\sum_{i=1}^{n}x_i-n/2}{\sqrt{n/12}}
\]

均值为零，标准偏差为 1。

\begin{enumerate}[start=2]
\item
  写一段小程序产生 (a) 中定义的变量 \(z\)，\(n\) 可以为任意值。分别取
  \(n=1,\ldots,20\)，生成直方图，每个直方图包含 \(10^4\) 个 \(z\)
  的值。何时 \(z\)
  的分布近似为高斯分布？作为简单的高斯分布产生子，可以取
  \(n=12\)。评论此算法的局限性。选作：对于 \(n=2\)，推导出 \(z\)
  的概率密度函数的明显形式。
\end{enumerate}
\begin{solution}
设 \(x_i\sim U(0,1)\) 独立同分布，
\[
y=\sum_{i=1}^n x_i.
\]
因为 \(E[x_i]=1/2\)、\(V[x_i]=1/12\)，独立性给出
\[
E[y]=\sum_i E[x_i]=\frac n2,
\qquad
V[y]=\sum_i V[x_i]=\frac n{12}.
\]
标准化变量
\[
z=\frac{y-n/2}{\sqrt{n/12}}
\]
满足
\[
E[z]=0,
\qquad
V[z]=1.
\]
由中心极限定理，\(n\) 增大时 \(z\) 的分布趋近标准高斯分布。实际作图时，\(n=1\)
仍是均匀分布，\(n=2\) 是三角形分布，\(n\approx 10\sim12\) 时中心部分已经相当接近高斯。
\pyfig[0.78\linewidth]{figures/fig_03_04_clt.png}{不同 \(n\) 下标准化变量 \(z\) 的经验密度；\(n\) 增大时逐渐接近标准高斯。}
这个算法的局限性是：\(z\) 的取值范围有限，
\[
-\sqrt{3n}<z<\sqrt{3n},
\]
所以它不能正确产生高斯分布的极端尾部；例如 \(n=12\) 时只能产生 \([-6,6]\)
内的值。此外，尾部形状与真正高斯尾部仍有系统差异。

选作：当 \(n=2\) 时，\(y=x_1+x_2\) 的密度为
\[
f_Y(y)=
\begin{cases}
y,&0<y<1,\\
2-y,&1<y<2,\\
0,&\text{其它},
\end{cases}
\]
而 \(z=(y-1)/\sqrt{1/6}=\sqrt6(y-1)\)。故
\[
f_Z(z)=
\begin{cases}
\dfrac{1+z/\sqrt6}{\sqrt6},&-\sqrt6<z<0,\\[4pt]
\dfrac{1-z/\sqrt6}{\sqrt6},&0<z<\sqrt6,\\[4pt]
0,&\text{其它}.
\end{cases}
\]
\end{solution}
\end{question}


\begin{question}
变量 \(t\) 服从均值 \(\tau=1\) 的指数分布，\(x\) 服从均值
\(\mu=0\)、标准偏差 \(\sigma=0.5\) 的高斯分布。写一段 MC 程序，产生变量

\[
y=t+x.
\]

这里 \(t\) 可以表示不稳定粒子的真实衰变时间，\(x\) 表示测量误差，所以
\(y\) 表示测量到的衰变时间。画出 \(y\) 的直方图。修改程序以研究
\(\tau\ll\sigma\) 和 \(\tau\gg\sigma\) 的情形。
\begin{solution}
设 \(t\sim\operatorname{Exp}(\tau)\)，即 \(E[t]=\tau\)、\(V[t]=\tau^2\)，
且 \(x\sim N(0,\sigma^2)\)，二者独立。则
\[
y=t+x
\]
服从指数分布与高斯分布的卷积，常称为 exGaussian 分布。基本矩为
\[
E[y]=E[t]+E[x]=\tau,
\qquad
V[y]=V[t]+V[x]=\tau^2+\sigma^2.
\]
若需要写出解析密度，则
\[
f_Y(y)=\frac{1}{2\tau}
\exp\left(\frac{\sigma^2}{2\tau^2}-\frac{y}{\tau}\right)
\operatorname{erfc}\left(\frac{\sigma^2/\tau-y}{\sqrt2\sigma}\right).
\]

MC 程序中可先产生 \(t=-\tau\log r\)（\(r\sim U(0,1)\)），再产生一个高斯随机数
\(x\)，最后填入 \(y=t+x\) 的直方图。
\pyfig[0.76\linewidth]{figures/fig_03_05_exgaussian.png}{不同 \(\tau/\sigma\) 比例下 \(y=t+x\) 的模拟分布。}
\begin{enumerate}
\item 当 \(\tau\ll\sigma\) 时，测量误差主导，分布中心部分近似高斯，只带很弱的右尾。
\item 当 \(\tau\sim\sigma\) 时，分布明显右偏，既有高斯展宽又有指数长尾。
\item 当 \(\tau\gg\sigma\) 时，指数衰变形状主导，高斯误差主要抹平 \(y\) 接近零的边界。
\end{enumerate}
\end{solution}
\end{question}


\begin{question}
考虑服从 Cauchy(Breit-Wigner) 分布的随机变量 \(x\)，

\[
f(x)=\frac{1}{\pi}\frac{1}{1+x^2}.
\]

\begin{enumerate}
\item
  证明如果 \(r\) 均匀分布于 \([0,1]\)，则
\end{enumerate}

\[
x(r)=\tan\left[\pi\left(r-\frac{1}{2}\right)\right]
\]

服从 Cauchy 分布。

\begin{enumerate}[start=2]
\item
  利用 (a) 中的结果，写一段小程序产生 Cauchy 分布的随机数。产生 \(10^4\)
  个事例并画出直方图。
\item
  修改 (b) 中的程序，重复进行实验，每次实验 \(n\)
  个独立的柯西分布数值（如取 \(n=10\)）。对每个样本，计算样本均值
  \(\bar{x}=\frac1n\sum_{i=1}^{n}x_i\)。比较 \(\bar{x}\) 的直方图与
  \(x\) 的原始直方图。（参见第 10 章第 8 题。）
\end{enumerate}
\begin{solution}
\begin{enumerate}
\item 若 \(r\sim U(0,1)\)，令
\[
x=\tan\left[\pi\left(r-\frac12\right)\right].
\]
则
\[
r=\frac12+\frac1\pi\arctan x,
\]
所以分布函数为
\[
F(x)=P(X\le x)=\frac12+\frac1\pi\arctan x.
\]
对 \(x\) 求导，得到
\[
f(x)=F'(x)=\frac1{\pi(1+x^2)},
\]
即标准 Cauchy 分布。

\item 程序中只需对每个均匀随机数 \(r\) 计算上式的 \(x(r)\) 并填入直方图。由于 Cauchy
分布有很长的尾部，作图时常把横轴限制在例如 \([-10,10]\) 或 \([-20,20]\)，否则少数极端值会压缩中心区域。

\item 对每次实验产生 \(n\) 个独立 Cauchy 变量并计算样本均值
\[
\bar x=\frac1n\sum_{i=1}^n x_i.
\]
由第 10 章的特征函数结论，\(\bar x\) 仍服从同一个标准 Cauchy 分布；因此其直方图不会像高斯样本均值那样随 \(n\) 增大而变窄。这不是数值程序错误，而是因为 Cauchy 分布没有有限的均值和方差，通常的大数定律条件不成立。
\pyfig[0.76\linewidth]{figures/fig_03_06_cauchy_mean.png}{单个 Cauchy 变量与 10 个 Cauchy 变量样本均值的直方图；两者服从同一分布。}\end{enumerate}
\end{solution}
\end{question}


\begin{question}
光电倍增管是可以探测单光子的设备，如图 3.1 所示。
\pyfig[0.86\linewidth]{figures/fig_03_07_pmt_schematic.png}{图 3.1：光电倍增管级联放大示意图。}


光子打到光阴极上，有一定的概率发射出一个电子（即光电子）。光电子在电场中向电极（称为打拿极）加速。当光电子撞到第一个打拿极时，光电子可以进一步释放出电子。这些电子又被加速到第二个打拿极，产生更多的电子。该过程一直持续，直到最后一个打拿极产生的电子被收集起来。

进入第 \(i\) 个打拿极的每个电子产生的次级电子数目可以看成均值为
\(\nu_i\) 的泊松变量 \(n_i\)，一般而言，每个打拿极的泊松均值 \(\nu_i\)
不同。假设光电倍增管有 \(N\) 个打拿极，单个入射光电子最终产生的电子数目
\(n_{\rm out}\) 的期待值（即光电倍增管的增益）为

\[
\nu_{\rm out}=E[n_{\rm out}]=\prod_{i=1}^{N}\nu_i.
\]

\begin{enumerate}
\item 写一段蒙特卡罗程序，用来计算 \(N=6\) 时单个光电子产生的电子数
\(n_{\rm out}\) 的分布，对所有打拿极取 \(\nu=3.0\)。（ROOT
中可以直接调用函数产生泊松分布。）利用你写的程序，重复 \(M=1000\)
次单个光电子事件，从而得到 \(M\) 个 \(n_{\rm out}\)，作出这 \(M\) 个
\(n_{\rm out}\) 值的直方图。（建议直方图分 50
个区间。）通过计算样本均值和样本方差

\[
\bar n_{\rm out}=\frac{1}{M}\sum_{i=1}^{M}n_{{\rm out},i}
\]

\[
s_{\rm out}^2=\frac{1}{M-1}\sum_{i=1}^{M}(n_{{\rm out},i}-\bar n_{\rm out})^2
\]

以估计均值 \(\nu_{\rm out}\) 和方差
\(V[n_{\rm out}]=\sigma_{\rm out}^2\)（样本均值和样本方差详见
Statistical Data Analysis 第 5
章）。将样本均值与（3.7）式给出的值进行比较。比较样本方差（或标准差）与均值为
\(\nu_{\rm out}\) 的泊松变量的方差。定性解释为什么 \(n_{\rm out}\)
的标准差远大于泊松变量的标准差。

\item 我们希望 \(n_{\rm out}\)
的标准差尽量小，以便能够尽可能精确地确定初始光电子的数目（从而可以估计入射光子的数目）。在某些应用中，需要标准差小到可以区分到底是一个还是两个光子，因此希望使相对分辨率（即标准差与均值的比值）小于
1。减小方差的一个方法是提高第一个打拿极产生电子数的均值，这可以通过增大电压从而提高入射到打拿极的光电子的能量来实现，也可以通过选取较低功函数的金属（即提高发射次级电子的概率）来实现。

修改（a）中的程序，增大第一个打拿极的均值（比如增大至
6）。运行程序并估计 \(n_{\rm out}\)
的标准差与均值的比值。定性解释为什么这么做比所有 \(\nu_i\)
都相同时的分辨率更好。为什么提高后面的打拿极的增益对提高分辨率作用不大？

\item 尝试将程序改成模拟 \(N=12\)
个打拿极。你会发现模拟每个打拿极上每个电子的碰撞需要太长时间。可以考虑换个思路，取
\(N=6\) 运行足够多的事例从而获得足够精确的 \(n_{\rm out}\)
的分布（例如至少 \(M\approx10^4\)，分 50 个区间作出
\(0\le n_{\rm out}\le5000\)
之间的直方图）。用某种方法，比如舍选法，产生服从该分布的随机数。对前 6
个打拿极得到的每个电子，用同样方法模拟它在后面 6
个打拿极产生的电子数目。对前 6 个打拿极，取 \(\nu_1=6\)，其余
\(\nu_i=3\)；对后 6 个打拿极，全部取 \(\nu_i=3\)。
\end{enumerate}
\begin{solution}
令 \(Z_i\) 表示第 \(i\) 个打拿极之后的电子数，初始 \(Z_0=1\)。在给定
\(Z_{i-1}\) 时，若每个电子在第 \(i\) 级产生均值为 \(\nu_i\) 的泊松个数，则
\[
Z_i\mid Z_{i-1}\sim\operatorname{Pois}(\nu_i Z_{i-1}).
\]
于是
\[
E[Z_i]=\nu_iE[Z_{i-1}],
\qquad
V[Z_i]=\nu_iE[Z_{i-1}]+\nu_i^2V[Z_{i-1}].
\]
递推得到
\[
E[n_{\rm out}]=\prod_{i=1}^N\nu_i,
\]
以及
\[
\frac{V[n_{\rm out}]}{E[n_{\rm out}]^2}
=\sum_{k=1}^N\frac{1}{\prod_{j=1}^k\nu_j}.
\]
\pyfig[0.76\linewidth]{figures/fig_03_07_pmt_hist.png}{六级 PMT 的输出电子数模拟；提高第一级平均增益会显著改善相对分辨率。}
\begin{enumerate}
\item 当 \(N=6\)、所有 \(\nu_i=3\) 时，
\[
E[n_{\rm out}]=3^6=729,
\]
\[
\sigma_{\rm out}=729\sqrt{\frac{1-3^{-6}}{3-1}}\simeq515.
\]
若把输出误认为均值 729 的单个泊松变量，则标准差只有 \(\sqrt{729}=27\)。真实标准差大得多，是因为前几级的涨落会被后续级联倍增并放大。

\item 若把第一级提高到 \(\nu_1=6\)，其余仍为 3，则均值变为
\(6\cdot3^5=1458\)，相对分辨率为
\[
\frac{\sigma}{E}
=\sqrt{\frac16+\frac1{18}+\frac1{54}+\frac1{162}+\frac1{486}+\frac1{1458}}
\simeq0.500.
\]
而所有级均为 3 时相对分辨率约为 \(0.707\)。提高第一级最有效，因为第一级涨落会被所有后续级放大；提高最后几级主要改变总增益，对已经形成的相对涨落改善很小。

\item 直接逐个电子模拟 12 个打拿极会非常慢。更好的做法是利用条件分布
\[
Z_i\mid Z_{i-1}\sim\operatorname{Pois}(\nu_i Z_{i-1})
\]
逐级抽样一次，而不是对每个电子逐个循环。对
\((\nu_1,\ldots,\nu_{12})=(6,3,3,\ldots,3)\)，理论均值为
\[
E[n_{\rm out}]=6\cdot3^{11}=1062882,
\]
标准差由上面的递推公式给出，约为
\[
\sigma_{\rm out}\simeq5.31\times10^5.
\]
这种逐级泊松抽样与逐个电子的完整级联模拟等价，但计算量小得多。
\end{enumerate}
\end{solution}
\end{question}

